% Plantilla latex para publicación de artículo científico en la revista "La Mecatrónica en México"
%
% @author {Juan Manuel Aviña Muñoz}
% @email {juan.avina.m@gmail.com}
% @year 2023
% @version 1.0
% @license CC BY 4.0
%
% Se agradecen sus comentarios, contribuciones, reporte de bugs y la difusión de esta plantilla.




% Tipo de documento ytamaño de letra
\documentclass[10pt]{article}

% Ajuste de lenguaje
\usepackage[spanish, english]{babel}

% Tamaño de página y margenes
\usepackage[letterpaper, top=2.5cm,bottom=3.5cm, left=3cm, right=2.5cm, marginparwidth=1.75cm]{geometry}

% Permite manejo de ecuaciones y símbolos matemáticos
\usepackage{amsmath, amssymb, amsthm, amsfonts}

% Paquete para agregar imágenes
\usepackage{graphicx}

% Paquete para hipervinculos
\usepackage{hyperref}
\hypersetup{
    colorlinks=true,
    linkcolor=black,
    urlcolor=black,
    citecolor=black}

% Package for specifying fonts
\usepackage{fontspec}

% Permite el manejo de caracteres especiales en el español
\usepackage[utf8]{inputenc} 
\usepackage[T1]{fontenc}

% Subrayado de textos
\usepackage[normalem]{ulem}

% Personalizar los encabezados y pies de página
\usepackage{fancyhdr}
\pagestyle{fancy}

% Ruta para imagenes
\graphicspath{{image/}}

% Fuente principal del texto: Arial
\setmainfont{Arial}

% Ajusta la sangría
\setlength{\parindent}{1cm}

% Ajusta el espacio entre renglones
\setlength{\parskip}{0.5cm}

% Headers config
% Cambiar Volumen, Numero y páginas
\lhead{
    La Mecatrónica en México, Enero 20xx, Vol. x, No. x, páginas xx – xx\\Disponible en línea en \href{http://www.mecamex.net/revistas/LMEM}{www.mecamex.net/revistas/LMEM}\\ISSN: 2448-7031, Asociación Mexicana de Mecatrónica A.C}
\rhead{
    \includegraphics[width=1.2cm]{logo}}
\headsep=60pt


% Información del artículo
\title{\Huge Titulo de investigación\\(no mayor a dos renglones)}
\author{\Large Autores de la publicación}
\date{\normalsize Datos de la Institución de adscripción de los autores.\\
\textit{Especificar correo electrónico del autor}}



\begin{document}
\maketitle


\selectlanguage{spanish}
\begin{abstract}
    \textit{\normalsize El resumen describe de forma concisa y clara el contenido del artículo. Las ideas fundamentales de relevancia, así como los resultados y las conclusiones generales del trabajo reportado.}
    
    \vspace{0.5cm}

    \textbf{Palabras clave:} Palabras que por sí mismas dan idea del tema o área del artículo.
\end{abstract}


\selectlanguage{english}
\begin{abstract}
    \textit{\normalsize The abstract concisely and clearly describes the content of the article. The fundamental ideas of relevance, as well as the results and general conclusions of the reported work.}
    
    \vspace{0.5cm}

    \textbf{Keywords:} Words that by themselves give an idea of the subject or area of the article.
\end{abstract}



\selectlanguage{spanish}
\begin{center}
    \section{Introducción}
\end{center}

    La introducción es el texto que sirve para ir induciendo al lector hacia el tema que trata el artículo. Las citas a otras fuentes han de referenciarse con paréntesis cuadrados, como se muestra aquí \cite{greenwade93}. La numeración de dichas referencias se efectuará de acuerdo al orden de aparición que estas tengan a lo largo del documento.

    Generalmente, la introducción finaliza con alguna información relacionada al artículo o con algunas reflexiones sobre el resultado del trabajo que se está presentando. Al terminar una sección se dejarán dos renglones en blanco, antes de colocar el título de la siguiente sección.

    Se pide a los autores que los artículos abarquen un mínimo de 10 páginas. No se considera un límite máximo eso depende del trabajo realizado y lo que se desea difundir.



\newpage % Agrega o elimina estos saltos de página de acuerdo a tus necesidades.
\vspace{0.5cm}
\begin{center}
    \section{Título}
\end{center}

    El nombre de éste segundo título dependerá de cómo se haya estructurado el artículo.
    
    Se debe considerar una sección en la cual se presente la metodología a seguir, utilizando de preferencia un mapa conceptual (esquemas, diagrama a bloques, de flujo, etc.) para mostrar gráficamente la metodología utilizada. Algunas de las secciones que pudieran llevar la estructura del artículo son: Descripción del problema, Metodología, Modelación y Simulación, Diseño, Manufactura, Integración, Metodología, etc.

    \subsection{Título de subsección}

        Cuando se muestre una ecuación en el artículo, dicha ecuación ha de enumerarse. La numeración se efectuará de acuerdo al orden de aparición de la ecuación en el artículo. Cuando se haga referencia a una ecuación debe hacerse de una manera similar a: “… como se observa en \ref{eq:pend}”, donde \ref{eq:pend} es el número de ecuación correspondiente. A continuación se muestra un ejemplo de una ecuación y su numeración.

        \vspace{0.5cm}
        \begin{equation}
            m=\frac{y_2-y_1}{x_2-x_1}
            \label{eq:pend}
        \end{equation}
        \vspace{0.5cm}

        Es importante mantener la calidad de la información de los artículos, por esta razón se sugiere al autor mostrar y describir las figuras que él considere importantes en su documento. Por ejemplo, en la figura \ref{fig:frog} se muestra una rana.

        De preferencia todas las figuras deben tener su referencia, excepto cuando la figura es producto del trabajo relacionado con el trabajo, o diseñado por los autores.
        
        \vspace{0.5cm}
        \begin{figure}[h!]
        \centering
            \includegraphics[width=0.5\textwidth]{frog}
            \caption{\texttt{\textsl{{\footnotesize Redeyed Tree Frog Smile. [Referencia]}}}}
            \label{fig:frog}
        \end{figure}
        \vspace{0.5cm}

        Se debe agregar un diagrama a bloques que ilustre la metodología utilizada en el desarrollo del trabajo, y posteriormente se debe detallar cada uno de los bloques.

        \subsubsection{Título de sub-subsección}

            En caso de que sea necesario agregar una sección de tercer nivel, sigue este formato.



\begin{center}
    \section{Resultados}
\end{center}

        Aquí se muestran los resultados obtenidos, se pueden adjuntar tablas, citándolas de la siguiente manera \ref{tab:1}.

        \vspace{0.5cm}
        \begin{table}[h!]
            \centering
                \begin{tabular}{|c|c|c|c|}
                \hline
                \textbf{Característica} & \textbf{1} & \textbf{2} & \textbf{3} \\ \hline
                {Tipo 1}         & x          & x          & x          \\ \hline
                {Tipo 2}         & x          & x          & x          \\ \hline
                {Tipo 3}         & x          & x          & x          \\ \hline
                {Tipo 4}         & x          & x          & x          \\ \hline
            \end{tabular}
            \label{tab:1}
        \end{table}
        \vspace{0.5cm}

        Las tablas se deberán presentar en la medida posible, completas en la misma hoja del texto que le acompaña.



\begin{center}
    \section{Conclusiones}
\end{center}

        En esta sección por lo general se describen los avances del trabajo, así como una reflexión con base en los resultados del trabajo realizado. Es común también mostrar ideas hacia trabajos futuros o aportaciones del trabajo realizado. Las conclusiones deben ser claras y concisas procurando enunciar los principales hallazgos o descubrimientos realizados.



\vspace{0.5cm}
% Para agregar o eliminar referencias, revisa el archivo "ref.bib", ahí encontrarás algunos ejemplos.
\nocite{*}
\bibliographystyle{IEEEtran}
\bibliography{ref}

\begin{center}
    \section*{Autores}
\end{center}

    \textbf{Autor 1.} En ésta sección se muestra una breve biografía académica del autor, en no más de 150 palabras. El nombre completo del autor de acuerdo a como se muestra al inicio del trabajo.

    \textbf{Autor 2.} En ésta sección se muestra una breve biografía académica del autor, en no más de 150 palabras. El nombre completo del autor de acuerdo a como se muestra al inicio del trabajo.

    \textbf{Autor 3.} En ésta sección se muestra una breve biografía académica del autor, en no más de 150 palabras. El nombre completo del autor de acuerdo a como se muestra al inicio del trabajo.
    
\end{document}